% !PRINTED 774
\weberpage{774}
Now, understanding by $t$ a variable, we set:

\[
\begin{aligned}
{}[t - \chi(1)][t - \chi(r)] \ldots [t - \chi(r^{n-1})] &= \\
t^n + A_1t^{n-1} + A_2t^{n-2} + \cdots + A_n, &
\end{aligned}
\]

in which the coefficients $A_1, A_2, \ldots, A_n$, being integral symmetric functions of the $n$ quantities $\chi(1), \chi(r), \ldots, \chi(r^{n-1})$, are integral rational numbers. If, however, we set in it

\[
t = \chi(r) = \dfrac{[A(\alpha)]^n}{p}
\]

then the left side vanishes, and multiplication by $p^n$ gives:

\weq{14}{[A(\alpha)]^{n^2} = p\,C(\alpha),}

where $C(\alpha)$ denotes a number of the form (9) with \gesperrt{integral coefficients}.

If formula (14) is repeatedly raised to the $p$th power and theorem (1) is applied, there follows from it, for every power that is not less than $n^2$:

\weq{15}{A(\alpha^{p^\lambda}) = p\,\xi(\alpha),}

in which the integral function $\xi(\alpha)$ likewise has integral coefficients.

Here we may assume $\lambda$ to be divisible by $\varphi(m)$, and then, since $m$ is not divisible by $p$, we have by the generalised theorem of \gesperrt{Fermat} (Bd.~II, \S.\,14):

\[
p^\lambda \equiv 1 \ (\mathrm{mod}\ m), \quad \alpha^{p^\lambda} = \alpha,
\]

so that from (14) there follows: % [sic] in the source; equation (16) follows mathematically from (15), not (14)

\weq{16}{A(\alpha) = p\,\xi(\alpha).}

Hence in (13) the coefficients $a_0, a_1, \ldots, a_{\nu-1}$ are all divisible by $p$, and $a$ is consequently not divisible by $p$. In the same way it is shown that $b_0, b_1, \ldots, b_{\nu-1}$ are divisible by $p$, while $b$ is indivisible by $p$.

Accordingly there follows from (12), if we set $A(\alpha)B(\alpha) = p^2\Theta(\alpha)$:

\weq{17}{ab = p\,\Theta(\alpha),}

where $\Theta(\alpha)$ denotes a function of $\alpha$ of the form (9) with integral coefficients. But equation (17), since it is at most of degree $(\nu-1)$, would have to be identical with respect to $\alpha$, and $ab$ would have to be divisible by $p$, which is not possible.
